\scanpage{140}
\mspage{56}

En utilisant (\unsure{dans} chacune de ces formules) des deux côtés les équivalences de catégories

\[
\left\{\begin{aligned}
M&\xrightarrow{\approx}\underline{\operatorname{Hom}}^{!}(M^{\unsure{\mathrm{o}|b}},\mathrm{Ens})\\
N&\xrightarrow{\approx}\underline{\operatorname{Hom}}^{!}(N^{\mathrm{o}},\mathrm{Ens})
\end{aligned}\right.
\tag*{(4.1.13)}
\]

On peut dire \unsure{aussi|qu'ici} que la catégorie $M\overline{\boxtimes}N$ est équivalente à celle des \uline{couples} de foncteurs

\[
\left\{\begin{aligned}
\varphi &: M^{\mathrm{o}}\longrightarrow N\\
\psi &: N^{\mathrm{o}}\unsure{\longrightarrow|\text{---}}M
\end{aligned}\right.
\tag*{(4.1.\unsure{14|13})}
\]

\deleted{avec $\varphi\simeq\unsure{\psi}$} \deleted{\unsure{avec deux}} \deleted{\unsure{et}} munis \unsure{de|des} deux morphismes d'adjonction

\[
\left\{\begin{aligned}
\psi\varphi^{\mathrm{o}}&\xleftarrow{\alpha}\mathrm{id}_{M}\\
\varphi\psi^{\mathrm{o}}&\xleftarrow{\beta}\mathrm{id}_{N}
\end{aligned}\right.
\tag*{(4.1.15)}
\]

satisfaisant les deux compatibilités habituelles, exprimant l'adjonction de $\varphi^{\mathrm{o}}$ et de $\psi$ (ou, de façon équivalente, de $\psi^{\mathrm{o}}$ et de $\varphi$). Le foncteur

\[
(\varphi,\psi)\longmapsto\varphi\qquad(\text{resp. }(\varphi,\psi)\longmapsto\psi)
\]

donne une équivalence de la catégorie de ces couples avec $\deleted{\unsure{M}}\underline{\operatorname{Hom}}^{!}(M^{\mathrm{o}},N)$ (resp avec $\underline{\operatorname{Hom}}^{!}(N^{\mathrm{o}},M)$), de sorte qu'on a un carré commutatif d'équivalences \unsure{canoniques} de foncteurs :
% archive sheet number: 139
